% !TeX root = ../../MainBook.tex
\section{拉普拉斯变换}\label{sec:Laplace_Transform}

十八世纪，欧拉（Euler）在研究微分方程时，首先使用了积分变换的方法，将微分方程转化为代数方程来求解，随后，拉格朗日（Lagrange）和拉普拉斯（Laplace）也使用了类似的方法。其中，拉普拉斯所使用的积分变换为$y_s=\int \e^{-sx}\phi(x)\,\dd x$，与现代的拉普拉斯变换最为相似，但当时并未以拉普拉斯变换命名。十九世纪，Petzval在其著作中首次系统地考察了这一方法，但他的学生Spitzer强烈主张将这一变换命名为拉普拉斯变换，这影响到几位很有影响力的数学家：Boole，Poincaré和Doetsch，他们在著作中均将这种方法称为“拉普拉斯变换”，使得这个名称被逐渐巩固。

%二十世纪初，Heaviside在研究电磁学时提出了一种符号运算的方法，将微分方程中的微分算子$\dfrac{d}{dt}$替换为符号$p$（或$D$），从而将微分方程转化为代数方程来求解。这种方法直观、快捷但缺乏严格的数学基础，随着数学家对拉普拉斯变换的研究不断深入，到1950年，拉普拉斯变换已完全取代Heaviside演算。

鉴于拉普拉斯变换的应用场景中使用角频率明显多于使用频率，我们在本书中采用角频率$\omega$作为拉普拉斯变换的变量；另外，本章默认$u(0)=1$。

\subsection{拉普拉斯变换的引出}\label{sub:introduction_of_laplace}
在\refsub{introduction_of_fourier}中，我们介绍了傅里叶变换：
\begin{equation}
    \mathcal{F}f(\omega)=\int_{-\infty}^{\infty}f(t)\e^{-\im\omega t}\,\dd t.
\end{equation}

在不考虑收敛性的问题时，可以将$\omega$推广为复数$s=\sigma+\im\omega$，根据含参变量积分的求导法则，这样的积分变换可以对$s$求导，从而是全纯函数。这个积分变换就是\textbf{双边拉普拉斯变换}：
\begin{definition}[双边拉普拉斯变换]
    \label{def:bi_laplace}
    设函数$f$局部可积，$\exists a,b\in\mathbb{R},\ a<b$（我们允许$a=-\infty$和$b=\infty$），使得$f(t)\e^{-at}$和$f(t)\e^{-bt}$均在$\mathbb{R}$上可积，$f$的双边拉普拉斯变换定义为
    \begin{equation}
        \mathcal{L}f(s)=\int_{-\infty}^{\infty}f(t)\e^{-st}\,\dd t,\quad s=\sigma+\im\omega,\quad a\le\sigma\le b.
    \end{equation}
    如果$f(t)$是时域信号，$\mathcal{L}f(s)$称为\textbf{复频域信号}。

    我们仍记$\mathcal{L}f(s)$为$F(s)$或$\hat{f}(s)$（在涉及傅里叶变换时，我们会避免使用它们或加以说明），并记$f\overset{\mathcal{L}}{\longleftrightarrow}F(s)$。
\end{definition}
“局部可积”即$f$在任意有限区间上可积，这个条件是为了保证在添加适当的收敛因子$\e^{-st}$后，拉普拉斯积分收敛，从而拉普拉斯变换有定义。

在$\sigma>0$时，这个收敛因子会使得$t\to-\infty$时$\abs{\e^{-st}}=\e^{-\sigma t}$快速增长，从而对被变换的函数$f$提出很高的要求。事实上，多数情况下$f$支于$[a,\infty),\ a\in\mathbb{R}$为常数，问题化归为因果信号（\cref{def:casual_signal}）的拉普拉斯变换，即\textbf{单边拉普拉斯变换}\cite{folland}：
\begin{definition}[单边拉普拉斯变换]
    \label{def:single_laplace}
    记
    \begin{equation}
        \mathcal{E}=\braces{f:[0,\infty)\to\mathbb{C}\;\middle|\;f\in PC[0,\infty),\ \exists C>0,\ a\in\mathbb{R},\ \abs{f(t)}<C\e^{at},\ \forall t\geq 0},
    \end{equation}
    其中$PC[0,\infty)$表示$[0,\infty)$上的分段连续函数。
    对$f\in\mathcal{E}$，其单边拉普拉斯变换定义为
    \begin{equation}
        \mathcal{L}f(s)=\int_{0}^{\infty}f(t)\e^{-st}\,\dd t,\quad s=\sigma+\im\omega,\quad \sigma>a.
    \end{equation}
\end{definition}
在$\mathcal{E}$的定义中我们对$f$的定义域做了要求，但实际上只要$f$是因果信号，就可以认为$f\in\mathcal{E}$，因此后文中不会特意区分$f$的定义域。

下面几节中，我们将介绍单边拉普拉斯变换的性质，它们大多容易推广到双边拉普拉斯变换的情形。如无特殊说明，本节后文中所说的拉普拉斯变换均指单边拉普拉斯变换。

在$\mathcal{E}$上定义拉普拉斯变换，是因为对于$f\in\mathcal{E}$，当$\sigma>a$时，给出拉普拉斯变换的积分是绝对收敛的：
\begin{equation}
    \int_0^{\infty}\abs{f(t)\e^{-st}}\,\dd t\leq \int_0^{\infty}C\e^{at}\e^{-\sigma t}\,\dd t=\frac{C}{\sigma-a}<\infty.
\end{equation}
因此在导数存在的情况下可以交换积分与求导的次序，从而$\mathcal{L}f(s)$是$\sigma>a$上的全纯函数：
\begin{align}
    \frac{\dd}{\dd s}\mathcal{L}f(s)&=\frac{\dd}{\dd s}\int_0^{\infty}f(t)\e^{-st}\,\dd t\nonumber\\
    &=\int_0^{\infty}f(t)\frac{\partial}{\partial s}(\e^{-st})\,\dd t\nonumber\\
    &=-\int_0^{\infty}tf(t)\e^{-st}\,\dd t\nonumber\\
    &=-\mathcal{L}[tf(t)](s).
\end{align}
这正是拉普拉斯变换的复频域微分性质。

根据全纯函数的唯一延拓性，如果$\mathcal{L}f(s)$可以延拓至$\sigma\leq a$的一些区域，则延拓是唯一的，尽管这时积分不一定存在，我们也认为延拓所得的定义域更大的函数是$f$的拉普拉斯变换。
\begin{example}[多项式的拉普拉斯变换]
    设$f(t)=t^\nu\ (\nu\geq 0)$，则$f\in\mathcal{E}$，且
    \begin{align}
        \mathcal{L}f(s)&=\int_0^{\infty}t^\nu \e^{-st}\,\dd t\nonumber\\
        &=\frac{1}{s^{\nu+1}}\int_0^{\infty}x^\nu \e^{-x}\,\dd x&(x=st)\nonumber\\
        &=\frac{\Gamma(\nu+1)}{s^{\nu+1}},\quad \sigma>0.
    \end{align}
    尽管积分只在$\sigma>0$时收敛，但$\mathcal{L}f(s)$可以延拓为$\mathbb{C}\backslash\{0\}$上的全纯函数。
    \begin{remark}
        其实，这个公式对实部大于$-1$的$\nu$都成立。尽管$-1<\Re(\nu)<0$时，$f(t)=t^\nu$在$t=0$处无界，但由于
        \begin{align}
            \abs{\int_{0}^{1}t^{\nu}\e^{-st}\,\dd t}&\leq \int_0^1 t^\nu\,\dd t=\frac{1}{\nu+1},\quad \sigma>0,\\
            \abs{\int_{1}^{\infty}t^{\nu}\e^{-st}\,\dd t}&\leq \int_1^\infty \e^{-\sigma t}\,\dd t=\frac{\e^{-\sigma}}{\sigma},\quad \sigma>0,
        \end{align}
        其拉普拉斯变换是存在的。因此，我们可以认为$t^\nu\in\mathcal{E},\ \Re(\nu)>-1$。
    \end{remark}
\end{example}

在$\mathcal{E}$中，卷积总是良定义的：
\begin{proposition}[$\mathcal{E}$中的卷积]
    设$f,g\in\mathcal{E}$，则它们的卷积为
    \begin{equation}
        (f*g)(t)=\begin{dcases}
            0, & \text{if}\ t<0,\\
            \int_0^t f(x)g(t-x)\,\dd x, & \text{if}\ t\geq 0.
        \end{dcases}
    \end{equation}
\end{proposition}
这是因为$x\geq 0,\ t-x\geq 0$时，$f(x),g(t-x)$才不恒为$0$，这相当于积分区间为$[0,t]$。

\subsection{单边拉普拉斯变换的运算性质}\label{sub:laplace_operation}

\begin{proposition}[拉普拉斯变换的性质]
    设$f,g\in\mathcal{E}$，$a,b\in\mathbb{C}$，则：
    \begin{itemize}
        \item 线性：在$f,g$的公共收敛域内，有
        \begin{equation}
            \mathcal{L}[af(t)+bg(t)](s)=a\mathcal{L}f(s)+b\mathcal{L}g(s);
        \end{equation}
        \item 时移：
        \begin{equation}
            \mathcal{L}[f(t-t_0)u(t)](s)=\e^{-st_0}\mathcal{L}f(s);
        \end{equation}
        \item 频移：
        \begin{equation}
            \mathcal{L}[\e^{at}f(t)](s)=\mathcal{L}f(s-a);
        \end{equation}
        \item 伸缩：设$a>0$，则
        \begin{equation}
            \mathcal{L}[f(at)](s)=\frac{1}{a}\mathcal{L}f\paren{\frac{s}{a}};
        \end{equation}
        \item 时域微分：设$f\in C^1[0,\infty),\ f'\in\mathcal{E}$，则
        \begin{equation}
            \mathcal{L}[f'(t)](s)=s\mathcal{L}f(s)-f(0),
        \end{equation}
        $f'(t)$在$0$处的定义为$f$在$0$处的右导数：
        \begin{equation}
            f'(0):=\lim_{h\to 0+}\frac{f(h)-f(0)}{h}.
        \end{equation}
        注意我们必须要求$f$是连续的，否则微积分基本定理需要分段使用，会使公式更加复杂。得到时域微分性质至少应该要求$f$是连续且分段光滑的；
        \item 复频域微分：
        \begin{equation}
            \mathcal{L}[tf(t)](s)=-\frac{\dd}{\dd s}\mathcal{L}f(s);
        \end{equation}
        \item 时域积分：
        \begin{equation}
            \mathcal{L}\sbr{\int_0^t f(x)\,\dd x}(s)=\frac{1}{s}\mathcal{L}f(s);
        \end{equation}
        \item 复频域积分：设$\displaystyle\frac{f(t)}{t}\in\mathcal{E}$，则
        \begin{equation}
            \mathcal{L}\sbr{\frac{f(t)}{t}}(s)=\int_s^{\infty}\mathcal{L}f(x)\,\dd x,
        \end{equation}
        此处积分路径为复平面上满足$\Im(s)=\omega$有界的任意从$s$到无穷远处的路径。
    \end{itemize}
\end{proposition}
\begin{remark}
    很明显，时域微分性质和复频域微分性质都可以对更高阶的导数进行推广，其中，时域微分性质推广为
    \begin{align}
        \mathcal{L}\sbr{f^{(n)}(t)}(s)&=s^n\mathcal{L}f(s)-\sum_{k=0}^{n-1}s^{n-1-k}f^{(k)}(0)\nonumber\\
        &=s^n\mathcal{L}f(s)-s^{n-1}f(0)-s^{n-2}f'(0)-\cdots-f^{(n-1)}(0).
    \end{align}
\end{remark}
\begin{proof}
    1.\ 线性性：直接得自积分的线性性。

    2.\ 时移：
    \begin{align}
        \mathcal{L}[f(t-t_0)u(t)](s)&=\int_0^{\infty}f(t-t_0)\e^{-st}\,\dd t\nonumber\\
        &=\int_{t_0}^{\infty}f(t-t_0)\e^{-st}\,\dd t\nonumber\\
        &=\e^{-st_0}\int_0^{\infty}f(x)\e^{-sx}\,\dd x&(x=t-t_0)\nonumber\\
        &=\e^{-st_0}\mathcal{L}f(s).
    \end{align}

    3.\ 频移：
    \begin{equation}
        \mathcal{L}[\e^{at}f(t)](s)=\int_0^{\infty}\e^{at}f(t)\e^{-st}\,\dd t=\int_0^{\infty}f(t)\e^{-(s-a)t}\,\dd t=\mathcal{L}f(s-a).
    \end{equation}

    4.\ 伸缩：
    \begin{align}
        \mathcal{L}[f(at)](s)&=\int_0^{\infty}f(at)\e^{-st}\,\dd t\nonumber\\
        &=\int_0^{\infty}f(x)\e^{-sx/a}\,\frac{\dd x}{a}&(x=at)\nonumber\\
        &=\frac{1}{a}\mathcal{L}f\paren{\frac{s}{a}}.
    \end{align}

    5.\ 时域微分：
    \begin{align}
        \mathcal{L}[f'(t)](s)&=\int_0^{\infty}f'(t)\e^{-st}\,\dd t\nonumber\\
        &=\evalat{\paren{f(t)\e^{-st}}}{0}{\infty}+\int_0^{\infty}sf(t)\e^{-st}\,\dd t\nonumber\\
        &=s\mathcal{L}f(s)-f(0).
    \end{align}

    6.\ 复频域微分：已经在上一小节证明过。

    7.\ 时域积分：利用时域微分性质，有
    \begin{equation}
        \mathcal{L}\sbr{\int_0^t f(x)\,\dd x}(s)=\frac{1}{s}\mathcal{L}\sbr{\frac{\dd}{\dd t}\paren{\int_0^t f(x)\,\dd x}}(s)=\frac{1}{s}\mathcal{L}f(s).
    \end{equation}

    8.\ 复频域积分：利用复频域的微分性质，有
    \begin{align}
        \frac{\dd}{\dd s}\mathcal{L}\sbr{\frac{f(t)}{t}}(s)&=-\mathcal{L}f(s),\\
        \mathcal{L}\sbr{\frac{f(t)}{t}}(s)&=\int_s^{\infty}\mathcal{L}f(x)\,\dd x.
    \end{align}
\end{proof}

\begin{theorem}[卷积定理]
    设$f,g\in\mathcal{E}$，则
    \begin{equation}
        \mathcal{L}[(f*g)(t)](s)=\mathcal{L}f(s)\cdot\mathcal{L}g(s).
    \end{equation}
\end{theorem}
\begin{proof}
    与前面傅里叶变换的卷积定理证明类似：
    \begin{align}
        \mathcal{L}[(f*g)(t)](s)&=\int_0^{\infty}(f*g)(t)\e^{-st}\,\dd t\nonumber\\
        &=\int_0^{\infty}\paren{\int_0^t f(x)g(t-x)\,\dd x}\e^{-st}\,\dd t\nonumber\\
        &=\int_0^{\infty}\int_0^{\infty}f(x)g(t-x)u(t-x)\e^{-st}\,\dd x\,\dd t\nonumber\\
        &=\int_0^{\infty}f(x)\paren{\int_x^{\infty}g(t-x)\e^{-st}\,\dd t}\,\dd x\nonumber\\
        &=\int_0^{\infty}f(x)\e^{-sx}\paren{\int_0^{\infty}g(y)\e^{-sy}\,\dd y}\,\dd x&(y=t-x)\nonumber\\
        &=\mathcal{L}f(s)\cdot\mathcal{L}g(s).
    \end{align}
\end{proof}

以上所有性质的证明都与傅里叶变换类似，实际上我们不需要做这种重复性的证明工作。给定$s=\sigma+\im\omega$，令$\tilde{f}(t)=f(t)\e^{-\sigma t}$，则$\mathcal{L}f(s)=\mathcal{F}\tilde{f}(\omega)$，只要$\tilde{f}\in L^1(\mathbb{R})$，就可以将傅里叶变换的性质直接应用到拉普拉斯变换上。由于我们在$\mathcal{E}$中要求了$f(t)$的分段连续性，在收敛域$\sigma>a$内，$\tilde{f}(t)$总是属于$L^1(\mathbb{R})$的。下面给出一个例子：
\begin{example}[时移性质]
    设$f\in\mathcal{E},\ t_0>0$，令$\tilde{f}(t)=f(t)\e^{-\sigma t}$，则
    \begin{align}
        \mathcal{L}[f(t-t_0)u(t-t_0)](s)&=\mathcal{F}\sbr{\tilde{f}(t-t_0)u(t-t_0)}(\omega)\nonumber\\
        &=\e^{-\im\omega t_0}\mathcal{F}\tilde{f}(\omega)\nonumber\\
        &=\e^{-st_0}\mathcal{L}f(s).
    \end{align}
\end{example}
\begin{example}[卷积定理]
    设$f,g\in\mathcal{E}$，令$\tilde{f}(t)=f(t)\e^{-\sigma t},\ \tilde{g}(t)=g(t)\e^{-\sigma t}$，注意到
    \begin{align}
        (\tilde{f}*\tilde{g})(t)&=\int_{-\infty}^{\infty}\tilde{f}(x)\tilde{g}(t-x)\,\dd x\nonumber\\
        &=\int_0^{t}f(x)\e^{-\sigma x}g(t-x)\e^{-\sigma(t-x)}\,\dd x\nonumber\\
        &=\e^{-\sigma t}\int_0^{t}f(x)g(t-x)\,\dd x\nonumber\\
        &=\e^{-\sigma t}(f*g)(t),
    \end{align}
    因此有
    \begin{equation}
        \mathcal{L}[f*g](s)=\mathcal{F}\sbr{f*g(t)\e^{-\sigma t}}(\omega)=\mathcal{F}[\tilde{f}*\tilde{g}](\omega)=\mathcal{F}\tilde{f}(\omega)\cdot\mathcal{F}\tilde{g}(\omega)=\mathcal{L}f(s)\cdot\mathcal{L}g(s).
    \end{equation}
\end{example}
其余的证明，请读者自行完成。

\subsection{常用函数的单边拉普拉斯变换}\label{sub:ordinary_laplace}

\begin{example}[指数多项式]
    设$f(t)=t^\nu \e^{at}u(t),\ \Re(\nu)>-1$，利用拉普拉斯变换的频移性质，有
    \begin{align}
        \mathcal{L}f(s)&=\mathcal{L}\sbr{t^\nu \e^{at}u(t)}(s)\nonumber\\
        &=\mathcal{L}\sbr{t^\nu u(t)}(s-a)\nonumber\\
        &=\frac{\Gamma(\nu+1)}{(s-a)^{\nu+1}},\quad \sigma>a.
    \end{align}
    根据有理函数的理论，任意真分式函数都可以表示为若干个此类函数的线性组合，因此可以说$f$是指数多项式等价于$\mathcal{L}f$是真分式函数。
\end{example}

\begin{example}[正弦、余弦函数]
    方法一：利用欧拉公式化归为指数函数：
    \begin{align}
        \mathcal{L}[\cos(\omega_0 t)](s)&=\mathcal{L}[\cos(\omega_0 t)u(t)](s)\nonumber\\
        &=\frac{1}{2}\mathcal{L}[\e^{\im\omega_0 t}u(t)](s)+\frac{1}{2}\mathcal{L}[\e^{-\im\omega_0 t}u(t)](s)\nonumber\\
        &=\frac{1}{2}\paren{\frac{1}{s-\im\omega_0}+\frac{1}{s+\im\omega_0}}\nonumber\\
        &=\frac{s}{s^2+\omega_0^2},\quad \sigma>0,
    \end{align}
    \begin{align}
        \mathcal{L}[\sin(\omega_0 t)](s)&=\mathcal{L}[\sin(\omega_0 t)u(t)](s)\nonumber\\
        &=\frac{1}{2\im}\mathcal{L}[\e^{\im\omega_0 t}u(t)](s)-\frac{1}{2\im}\mathcal{L}[\e^{-\im\omega_0 t}u(t)](s)\nonumber\\
        &=\frac{1}{2\im}\paren{\frac{1}{s-\im\omega_0}-\frac{1}{s+\im\omega_0}}\nonumber\\
        &=\frac{\omega_0}{s^2+\omega_0^2},\quad \sigma>0.
    \end{align}
    方法二：利用微分性质：
    \begin{align}
        -\omega_0^2\mathcal{L}[\cos(\omega_0 t)](s)&=\mathcal{L}[\cos''(\omega_0 t)](s)\nonumber\\
        &=s^2\mathcal{L}[\cos(\omega_0 t)](s)-s\cos(0)-\cos'(0)\nonumber\\
        &=s^2\mathcal{L}[\cos(\omega_0 t)](s)-s,
    \end{align}
    而$\cos(\omega_0 t)$具有常数级的增长速度（在\refsub{converge_abscissa}中，我们将定量描述“增长速度”这一概念），因此
    \begin{equation}
        \mathcal{L}[\cos(\omega_0 t)](s)=\frac{s}{s^2+\omega_0^2},\quad \sigma>0.
    \end{equation}
    同理，有
    \begin{equation}
        \mathcal{L}[\sin(\omega_0 t)](s)=\frac{\omega_0}{s^2+\omega_0^2},\quad \sigma>0.
    \end{equation}
\end{example}

\begin{example}[双曲函数]
    因为
    \begin{equation}
        \sinh(\omega_0 t)=\frac{\e^{\omega_0 t}-\e^{-\omega_0 t}}{2},\quad \cosh(\omega_0 t)=\frac{\e^{\omega_0 t}+\e^{-\omega_0 t}}{2},
    \end{equation}
    所以
    \begin{align}
        \mathcal{L}[\sinh(\omega_0 t)](s)&=\frac{1}{2}\mathcal{L}[\e^{\omega_0 t}](s)-\frac{1}{2}\mathcal{L}[\e^{-\omega_0 t}](s)\nonumber\\
        &=\frac{1}{2}\paren{\frac{1}{s-\omega_0}-\frac{1}{s+\omega_0}}\nonumber\\
        &=\frac{\omega_0}{s^2-\omega_0^2},\quad \sigma>\abs{\omega_0},
    \end{align}
    \begin{align}
        \mathcal{L}[\cosh(\omega_0 t)](s)&=\frac{1}{2}\mathcal{L}[\e^{\omega_0 t}](s)+\frac{1}{2}\mathcal{L}[\e^{-\omega_0 t}](s)\nonumber\\
        &=\frac{1}{2}\paren{\frac{1}{s-\omega_0}+\frac{1}{s+\omega_0}}\nonumber\\
        &=\frac{s}{s^2-\omega_0^2},\quad \sigma>\abs{\omega_0}.
    \end{align}
\end{example}

\begin{example}[采样函数]
    首先，因为$\displaystyle\frac{\sin(\omega_0 t)}{t}$有界，所以$\displaystyle\frac{\sin(\omega_0 t)}{t}\in\mathcal{E}$。

    结合正弦函数的拉普拉斯变换和拉普拉斯变换的复频域积分性质，有
    \begin{align}
        \mathcal{L}\sbr{\frac{\sin(\omega_0 t)}{t}}(s)&=\int_s^{\infty}\mathcal{L}[\sin(\omega_0 t)](x)\,\dd x\nonumber\\
        &=\int_s^{\infty}\frac{\omega_0}{x^2+\omega_0^2}\,\dd x\nonumber\\
        &=\evalat{\arctan\frac{x}{\omega_0}}{s}{\infty}\nonumber\\
        &=\frac{\pi}{2}-\arctan\frac{s}{\omega_0}\nonumber\\
        &=\arctan\paren{\frac{\omega_0}{s}},\quad \sigma>0.
    \end{align}
\end{example}

\begin{example}[正弦积分函数]
    $\mathrm{Si}$函数定义见\refsub{int_sin}。
    \begin{equation}
        \frac{\sin(\omega_0 t)}{t}\leq C\e^{at},\quad a>0\implies\mathrm{Si}(\omega_0 t)\leq\frac{C}{\omega_0 a}\e^{at}\implies\mathrm{Si}(\omega_0 t)\in\mathcal{E}.
    \end{equation}
    利用正弦函数的拉普拉斯变换和拉普拉斯变换的时域积分性质，有
    \begin{align}
        \mathcal{L}[\mathrm{Si}(\omega_0 t)](s)&=\mathcal{L}\sbr{\int_0^{\omega_0 t}\frac{\sin(x)}{x}\,\dd x}(s)\nonumber\\
        &=\mathcal{L}\sbr{\int_0^{t}\frac{\sin(\omega_0 x)}{x}\,\dd x}(s)\nonumber\\
        &=\frac{1}{s}\mathcal{L}\sbr{\frac{\sin(\omega_0 t)}{t}}(s)\nonumber\\
        &=\frac{1}{s}\arctan\paren{\frac{\omega_0}{s}},\quad \sigma>0.
    \end{align}
\end{example}

\subsection{*单边拉普拉斯变换的级数求法}\label{sub:laplace_by_series}
我们知道，多项式的拉普拉斯变换都可以通过
\begin{equation}
    \mathcal{L}[t^n](s)=\frac{\Gamma(n+1)}{s^{n+1}}
\end{equation}
来计算，那么对于一些可以展开为幂级数的函数，不难想到逐项求拉普拉斯变换，下面讨论这种操作的适用范围。在此之前，我们先介绍一个引理：
\begin{lemma}
    设$f\in\mathcal{E},\ \abs{f(t)}\leq C\e^{at}$，则
    \begin{enumerate}
        \item $\forall x>a,\ \mathcal{L}f(x+\im y)\to 0,\quad \abs{y}\to\infty$；
        \item $\forall y,\ \mathcal{L}f(x+\im y)\to 0,\quad x\to\infty$。
    \end{enumerate}
\end{lemma}
\begin{proof}
    1.\ 令$g(t)=f(t)\e^{-xt}$，则$\abs{g(t)}\leq C\e^{-(x-a)t}$，因此$g\in L^1(\mathbb{R})$，根据Riemann-Lebesgue引理，有
    \begin{equation}
        \mathcal{L}f(x+\im y)=\mathcal{F}g(y)\to 0,\quad \abs{y}\to\infty.
    \end{equation}

    2.\ 利用Heine归并原理，任取实部趋于无穷的序列$\{z_n\}_{n=1}^\infty$，满足$\mathcal{L}f(z_n)$存在（这要求$\Re(z)>a$），我们要证明序列$\mathcal{L}f(z_n)\to 0$。
    \begin{equation}
        \int_{0}^{\infty}\abs{f(t)\e^{-z_n t}}\,\dd t\leq \int_0^{\infty}C\e^{at}\e^{-\Re(z_n)t}\,\dd t=C\int_0^{\infty}\e^{-(\Re(z_n)-a)t}\,\dd t=\frac{C}{\Re(z_n)-a}<\infty.
    \end{equation}
    由控制收敛定理（\cref{thm:DCT}），取$\abs{f(t)\e^{-z_1 t}}$作为控制函数，得到
    \begin{equation}
        \lim_{n\to\infty}\mathcal{L}f(z_n)=\lim_{n\to\infty}\int_0^{\infty}f(t)\e^{-z_n t}\,\dd t=\int_0^{\infty}\lim_{n\to\infty}f(t)\e^{-z_n t}\,\dd t=0.
    \end{equation}
\end{proof}

\begin{theorem}[拉普拉斯变换的级数求法]
    设$f\in\mathcal{E}$，且在$[0,\infty)$上有幂级数展开
    \begin{equation}
        f(t)=\sum_{n=0}^{\infty}a_n t^n,
    \end{equation}
    如果存在$\omega_0>0$，使得级数
    \begin{equation}
        \sum_{n=0}^{\infty}\abs{a_n}\frac{\Gamma(n+1)}{\omega_0^{n}}
    \end{equation}
    收敛，则在$\sigma>\omega_0$时，级数
    \begin{equation}
        \sum_{n=0}^{\infty}a_n\frac{\Gamma(n+1)}{s^{n+1}}
    \end{equation}
    收敛且等于$\mathcal{L}f(s)$，即$f$的拉普拉斯变换可以通过对其幂级数求拉普拉斯变换得到。
\end{theorem}
\begin{proof}
    由条件知级数$\displaystyle\sum_{n=0}^{\infty}a_n\frac{\Gamma(n+1)}{s^{n+1}}$在$\sigma>\omega_0$时绝对收敛、内闭一致收敛，因此只需证明其和函数等于$\mathcal{L}f(s)$。拉普拉斯变换是一种积分变换，我们需要将求和与积分交换，可以考虑使用控制收敛定理。由于
    \begin{align}
        &\abs{a_n\Gamma(n+1)\omega_0^{-n}}\to 0,\quad n\to\infty\nonumber\\
        &\implies \abs{a_n}\leq C\frac{\omega_0^{n}}{\Gamma(n+1)}\nonumber\\
        &\implies \abs{f(t)}\leq \sum_{n=0}^{\infty}\abs{a_n}t^n\leq C\sum_{n=0}^{\infty}\frac{(\omega_0 t)^{n}}{\Gamma(n+1)}=C\e^{\omega_0 t}\nonumber\\
        &\implies \forall s,\ \Re(s)=\omega_0+\phi>\omega_0,\quad \abs{f(t)\e^{-st}}\leq C\e^{-\phi t},
    \end{align}
    控制收敛定理的条件成立，换序即证。
\end{proof}

一般地，对于幂次不是整数的情况，可以推广以上结论：
\begin{corollary}*
    设$f\in\mathcal{E}$，且在$[0,\infty)$上有广义幂级数展开
    \begin{equation}
        f(t)=\sum_{n=0}^{\infty}a_n t^{n+\alpha},
    \end{equation}
    其中$0\leq\alpha<1$为常数，如果存在$\omega_0>0$，使得级数
    \begin{equation}
        \sum_{n=0}^{\infty}\abs{a_n}\frac{\Gamma(n+\alpha+1)}{\omega_0^{n}}
    \end{equation}
    收敛，则在$\sigma>\omega_0$时，级数
    \begin{equation}
        \sum_{n=0}^{\infty}a_n\frac{\Gamma(n+\alpha+1)}{s^{n+\alpha+1}}
    \end{equation}
    收敛且等于$\mathcal{L}f(s)$，即$f$的拉普拉斯变换可以通过对其广义幂级数求拉普拉斯变换得到。
\end{corollary}
